\section{Problems and access models}\label{sec:models}
All vector norms are Euclidean and all matrix norms are spectral unless
specified otherwise. The adjoint is denoted by $^*$. A Hermitian matrix
$H\succ0$ is supplied with public bounds
\begin{equation}\label{eq:spectral-promise}
                 \kappa^{-1}I\preceq H\preceq I,\qquad \kappa\geq1.
\end{equation}
Here $\kappa$ is a certified upper bound on its condition number; statements
requiring equality say so explicitly. For a nonzero $b\in\C^N$, our main
target is
\begin{equation}\label{eq:target}
                            q(H,b)=b^*H^{-1}b.
\end{equation}
An estimate has relative error $\epsilon$ if
$|\widehat q-q|\leq\epsilon q$. If $H=\nabla^2F(x)$ and
$b=\nabla F(x)$, this is the squared Newton decrement at $x$. More
generally, $b$ may be a linear-objective perturbation of the barrier gradient.
A theorem about an arbitrary pair $(H,b)$ becomes an optimization theorem
only after specifying an objective, an interior point, and their input
oracles. We give such realizations separately.

For the optimization sections, a barrier $\phi$ is a $C^3$ convex
function on the relative interior of its feasible affine slice, expressed
in independent coordinates and with positive definite restricted Hessian.
Its local norm is $\|h\|_x^2=h^T\nabla^2\phi(x)h$, with dual norm
$\|g\|_{x,*}^2=g^T\nabla^2\phi(x)^{-1}g$.
A self-concordant barrier satisfies
$|D^3\phi(x)[h,h,h]|\leq2\|h\|_x^3$ and diverges at the boundary;
a parameter $\nu$ additionally bounds
$|D\phi(x)[h]|\leq\sqrt\nu\|h\|_x$.
Parameters stated below refer to the specified barrier and coordinates,
not to every formulation of the same feasible set.
For maximizing $c^Tx$, the barrier subproblem is
$F_\eta(x)=\phi(x)-\eta c^Tx$, $\eta>0$.
A central point minimizes $F_\eta$; an analytic center minimizes $\phi$.
The Newton direction is $-\nabla^2\phi(x)^{-1}\nabla F_\eta(x)$ and its
squared decrement is $\|\nabla F_\eta(x)\|_{x,*}^2$.
At a central point, changing the multiplier by $\Delta\eta$ gives a
predictor squared decrement $(\Delta\eta)^2\|c\|_{x,*}^2$.
Equality-constrained versions use the Hessian and gradient restricted to
the tangent space, or the explicitly stated KKT equations. These intrinsic
quantities need not have the condition number of an augmented KKT matrix.

\begin{definition}[Sparse access and vector sampling]\label{def:access}
Row-sparse access to a matrix supplies the nonzero count in each requested
row, the location of its $j$th nonzero, and its value, at unit query cost.
The ordering is fixed for each input. A known bound $d$ limits the count.
Column-sparse access has the analogous meaning. For Hermitian matrices,
column access follows from row access by conjugation.

For $v\ne0$, $\SQ(v)$ supplies coordinate queries, the norm $\norm v$, and
independent samples from $p_v(i)=|v_i|^2/\norm v^2$. Zero coordinates have
zero sampling probability. These are exact oracles in our main model.
Full matrix $\SQ$ additionally supplies row and column norms, the Frobenius
norm, within-row and within-column squared-magnitude samples, and samples
of rows or columns proportional to their squared norms. A zero row or
column is identified by its zero norm and has no conditional sampling
request. This stronger interface includes sparse locations when a lower
bound explicitly grants both interfaces.
\end{definition}

Specifying these oracles matters: constructing a sampling data structure
from an arbitrary explicit vector generally entails an initial pass over
that vector. We do not charge this preprocessing to an algorithm whose
input contract already grants $\SQ(v)$. Conversely, a lower-bound reduction
must simulate every granted oracle; it cannot ignore hidden information in
norms or sampling probabilities. Matrix $\SQ$ is a classical interface and
by itself does not grant a coherent quantum state-preparation unitary.

Our coherent statements specify a unitary $U$ whose designated top-left
block is $H/\alpha$, together with controlled calls to $U,U^*$ and a
specified preparation of $b/\norm b$. The normalization $\alpha$ and the
allowed unitary completion are part of that access model. An exact
sparse-value oracle, which returns an entry as a number, is a different
interface. A black-box block-encoding lower bound need not hold for it.

We count oracle invocations and ideal real or complex arithmetic operations.
Independent uniform random draws from $[0,1]$ and exact comparisons with
computed real probabilities are unit-cost primitives in this model.
In particular, the exact rejection samplers below use ideal randomness;
their arithmetic bounds do not count a finite-bit implementation of these
random choices.
An index is one word; its bit length is $\lceil\log_2N\rceil$. Thus
``dimension-independent'' means no dependence on $N$ in these counts,
not an $N$-independent bit cost or input-loading cost. The notation
$\widetilde O$ suppresses logarithmic factors only in the displayed
parameters. Our direct classical upper bounds below use ordinary $O$
notation, with unspecified absolute constants in polynomial degrees and
local-walk exponents. Unless stated otherwise, $0<\epsilon\leq1/2$ and
$0<\zeta<1/2$ denote accuracy and failure probability.
Unless a different contract is stated, every classical or quantum lower
bound requires success probability at least $2/3$ on every promised input
and charges the worst-case number of queries. Explicit confidence
parameters and distributional expected-cost contracts take precedence
over this convention.

Besides one scalar, we consider a solution-sampling output: the algorithm
provides coordinate access and samples exactly from $p_{\widetilde x}$ for
one fixed vector satisfying
$\norm{\widetilde x-A^{-1}b}\leq\epsilon\norm{A^{-1}b}$.
This is weaker than writing all coordinates, and stronger than an
unspecified approximate output distribution. It does not include an exact
norm oracle for $\widetilde x$. A lower bound using a rare event must also
specify whether unflagged algorithmic failure may contaminate that event.
These distinctions prevent transferring bounds between incompatible output
contracts.

Finally, a public interval $\mu I\preceq H\preceq LI$ reduces to
\eqref{eq:spectral-promise} by using $H/L$ with $\kappa=L/\mu$.
The resulting inverse quadratic form must be divided by $L$.
